\subsection{一般行列式零境界とtwist後の中心指標}
\label{subsec:cubic-as4-boundary-R}
まず、AS正則性に依存しない中心の制限を記す。
ここで右Zhang twistは $x\star y=x\psi^{\deg x}(y)$ の規約を使う。

\begin{lemma}\label{lem:cubic-twist-central-weights}
$D$ を局所有限な連結次数付き整域、$\psi$ を次数付き自己同型とし、$A=D^\psi$ とする。
$A$ が同じ正次数 $d$ の非零中心元 $q_1,q_2$ を持ち、
$\psi(q_j)=\chi_jq_j$ であるとする。
$Z(D)_n\ne0$ なら $(\chi_1/\chi_2)^n=1$ である。
特に二指標の比が $-1$ なら $D$ の全奇数中心は零となる。
\end{lemma}
\begin{proof}
$\psi$ は有限次元空間 $Z(D)_n$ を保つので、非零固有ベクトル $F$、
$\psi(F)=\rho F$ を選べる。
$D$ での $F$ の中心性から、$A$ では
\[
 F\star v=\theta(v)\star F,\qquad
 \theta|_{D_1}=\rho^{-1}\psi^n|_{D_1}
\]
となる。一次生成を用いる場合の表示で書いたが、一般には
$\theta(x)=\rho^{-\deg x}\psi^n(x)$ で同じ計算が成立する。
$A$ も整域なので、$F$ のnormalizerは $Z(A)$ を各元ごとに固定する。
従って $\rho^{-d}\chi_j^nq_j=q_j$、すなわち
$\rho^d=\chi_j^n$ である。二つを比較して結論を得る。
\end{proof}

\begin{proposition}\label{prop:cubic-boundary-R}
\cite{AS4lowpinned}のFG2にある全 $R^\eta_\lambda$、
$\lambda\ne0,-1$, $\eta=\pm1$ はH3包絡にならない。
より強く、この表示の整域実現の全奇数中心は零である。
\end{proposition}
\begin{proof}
正符号の表示は
\[
\begin{gathered}
 x_0=x_1=x_2=x_5=1,\quad x_9=(\lambda^2-1)/16,\\
 x_{10}=-\frac{i\lambda^2}{4(\lambda+1)},\quad
 x_{16}=-\frac{i}{4(\lambda+1)},\quad
 x_{11}=x_{15}=\frac{i\lambda}{4(\lambda+1)},\\
 x_{14}=\frac{\lambda^4}{16(\lambda+1)^2},\quad
 x_{19}=\frac{\lambda^2}{16(\lambda+1)^2},\quad
 x_{20}=\frac1{16(\lambda+1)^2}
\end{gathered}
\]
であり、他の $x_j$ は零である。
$k=\lambda/(4(\lambda+1))$ とし、元の生成空間で
\[
 \phi=\begin{pmatrix}
 0&0&0&1\\0&0&\lambda k&0\\0&k/\lambda&0&0\\-k^2&0&0&0
 \end{pmatrix},\qquad
 C=(a+ikd,\ a-ikd,\ \lambda b+c,\ -\lambda b+c)
\]
を取る。$C$ は列を新基底とする行列を表し、
\[
 \det C=-4ik\lambda\ne0,\qquad
 C^{-1}\phi C=k\sigma,\quad \sigma=\operatorname{diag}(i,-i,1,-1).
\]
新基底を再び $a,b,c,d$ と書く。
六関係の $(1\otimes\phi)$ による像を $C$ 基底で書くと、
$t=(\lambda-1)/(\lambda+1)$ に対する次の代数 $A_t$ となる。
\[
\begin{aligned}
 a^2+b^2+\frac{t}{16}(cd+dc)&=0,&
 ab+ba-\frac{c^2+d^2}{16}&=0,\\
 ac+ca&=0,&ad-da&=0,\\
 bc+cb-2t da&=0,&bd+db&=0.
\end{aligned}
\]
この全六関係の一致を $\mathbb Q(i)(\lambda)$ 上で独立確認した。
比較にはこの六関係の固定した六係数がidentityとなる左逆を使っており、
$\lambda-1$ 等を割って特殊値を落としていない。
$\sigma$ はこの関係空間を保つことも直接確認した。
関係空間の等式 $R_{A_t}=(1\otimes\sigma)R_D$ から、
$\sigma$ は $D$ の関係も保つ。従って本稿の右twist規約では、$D$ を $R^+_\lambda$ の $C$ 基底表示とすると
\[
 A_t=D^{\sigma^{-1}}.
\]
ここでのtwistは中心を制限する補助であり、intrinsic同型の同定には使わない。

$A_t$ において $a^2,d^2$ は非零中心二次元である。
これをAS4存在定理から輸入する必要はない。
例えば
\[
 [a^2,b]=[a,ab+ba]=\tfrac1{16}[a,c^2+d^2]=0
\]
であり、他の $a^2$ の交換子も二つの二項関係から従う。
また
\[
 [c,ab+ba]=2t[da,a]=0
\]
なので、第二関係から $[c,d^2]=0$ を得る。
$d^2$ と $a,b$ の可換性も明らかである。
これら八つの三次恒等式には $\mathbb Q[t]$ 上の関係moduleの証明書があり、
$t=0$ を含む全値で成立する。
さらに $aa-bb$ と $cc-dd$ の係数汎関数がそれぞれ全六関係を消すので、
$a^2,d^2$ は非零である。

$\psi=\sigma^{-1}$ はこの二中心元に $-1,+1$ で作用する。
$D$ が整域なら補題\ref{lem:cubic-twist-central-weights}により全奇数中心は零となる。
元の $R A R^{\mathsf T}$ の非零反巡回性も全 $\lambda$ で検証したので、
AS4実現では $\mu=I$、従ってH3切断は中心三次を必要として矛盾する。
負符号は、形式パラメータ $\lambda$ を固定した全恒等式の $i\mapsto-i$ 共役で同じ結論を得る。
このため $\lambda=1$ や $\lambda=\pm i$ も欠落しない。
\end{proof}

六関係の輸送・可逆基底・最高tensorと、八つの分母なし中心恒等式は
\texttt{verify\_\allowbreak{}cubic\_\allowbreak{}as4\_\allowbreak{}four\_\allowbreak{}simple\_\allowbreak{}boundary\_\allowbreak{}R.py}
で再現する。FG2/FG6の具体表示のみを参照し、
全CY被覆・AS4十分性・通常同型分類は輸入していない。
